%
\begin{isabellebody}%
\setisabellecontext{Primitive{\isacharunderscore}{\kern0pt}Recursion}%
%
\begin{isamarkuptext}%
\section*{Primitive Recursion on Algebraic Data Types}%
\end{isamarkuptext}\isamarkuptrue%
%
\isadelimtheory
%
\endisadelimtheory
%
\isatagtheory
\isacommand{theory}\isamarkupfalse%
\ Primitive{\isacharunderscore}{\kern0pt}Recursion\isanewline
\ \ \isakeyword{imports}\ Main\isanewline
\isakeyword{begin}%
\endisatagtheory
{\isafoldtheory}%
%
\isadelimtheory
%
\endisadelimtheory
%
\begin{isamarkuptext}%
We can define functions recursively on algebraic data types in a way that is analogous to 
      previous recursive functions we discussed on natural numbers or lists. Below is a function
      that recursively computes the set of elements in a tree, and the list of elements in the
      tree traversed in order.%
\end{isamarkuptext}\isamarkuptrue%
\isanewline
\isacommand{fun}\isamarkupfalse%
\ tree{\isacharunderscore}{\kern0pt}set{\isacharcolon}{\kern0pt}{\isacharcolon}{\kern0pt}{\isachardoublequoteopen}{\isacharprime}{\kern0pt}a\ tree\ {\isasymRightarrow}\ {\isacharprime}{\kern0pt}a\ set{\isachardoublequoteclose}\ \isakeyword{where}\isanewline
{\isachardoublequoteopen}tree{\isacharunderscore}{\kern0pt}set\ Leaf\ {\isacharequal}{\kern0pt}\ {\isacharbraceleft}{\kern0pt}{\isacharbraceright}{\kern0pt}{\isachardoublequoteclose}\isanewline
{\isacharbar}{\kern0pt}\ {\isachardoublequoteopen}tree{\isacharunderscore}{\kern0pt}set\ {\isacharparenleft}{\kern0pt}Node\ l\ a\ r{\isacharparenright}{\kern0pt}\ {\isacharequal}{\kern0pt}\ insert\ a\ {\isacharparenleft}{\kern0pt}{\isacharparenleft}{\kern0pt}tree{\isacharunderscore}{\kern0pt}set\ l{\isacharparenright}{\kern0pt}\ {\isasymunion}\ {\isacharparenleft}{\kern0pt}tree{\isacharunderscore}{\kern0pt}set\ r{\isacharparenright}{\kern0pt}{\isacharparenright}{\kern0pt}{\isachardoublequoteclose}\isanewline
\isanewline
\isacommand{definition}\isamarkupfalse%
\ tree{\isacharunderscore}{\kern0pt}eg\ \isakeyword{where}\ {\isachardoublequoteopen}tree{\isacharunderscore}{\kern0pt}eg\ {\isasymequiv}\ Node\ {\isacharparenleft}{\kern0pt}Node\ \ Leaf\ {\isadigit{1}}\ {\isacharparenleft}{\kern0pt}Node\ Leaf\ {\isadigit{2}}\ Leaf{\isacharparenright}{\kern0pt}{\isacharparenright}{\kern0pt}\ {\isacharparenleft}{\kern0pt}{\isadigit{4}}{\isacharcolon}{\kern0pt}{\isacharcolon}{\kern0pt}nat{\isacharparenright}{\kern0pt}\ {\isacharparenleft}{\kern0pt}Node\ Leaf\ {\isadigit{3}}\ Leaf{\isacharparenright}{\kern0pt}{\isachardoublequoteclose}\isanewline
\isanewline
\isacommand{value}\isamarkupfalse%
\ {\isachardoublequoteopen}tree{\isacharunderscore}{\kern0pt}set\ tree{\isacharunderscore}{\kern0pt}eg{\isachardoublequoteclose}\isanewline
\isanewline
\isacommand{fun}\isamarkupfalse%
\ inorder{\isacharcolon}{\kern0pt}{\isacharcolon}{\kern0pt}{\isachardoublequoteopen}{\isacharprime}{\kern0pt}a\ tree\ {\isasymRightarrow}\ {\isacharprime}{\kern0pt}a\ list{\isachardoublequoteclose}\ \isakeyword{where}\isanewline
\ \ {\isachardoublequoteopen}inorder\ Leaf\ {\isacharequal}{\kern0pt}\ {\isacharbrackleft}{\kern0pt}{\isacharbrackright}{\kern0pt}{\isachardoublequoteclose}\isanewline
{\isacharbar}{\kern0pt}\ {\isachardoublequoteopen}inorder\ {\isacharparenleft}{\kern0pt}Node\ l\ a\ r{\isacharparenright}{\kern0pt}\ {\isacharequal}{\kern0pt}\ {\isacharparenleft}{\kern0pt}inorder\ l{\isacharparenright}{\kern0pt}\ {\isacharat}{\kern0pt}\ {\isacharbrackleft}{\kern0pt}a{\isacharbrackright}{\kern0pt}\ {\isacharat}{\kern0pt}\ {\isacharparenleft}{\kern0pt}inorder\ r{\isacharparenright}{\kern0pt}{\isachardoublequoteclose}\isanewline
\isanewline
\isacommand{value}\isamarkupfalse%
\ {\isachardoublequoteopen}inorder\ tree{\isacharunderscore}{\kern0pt}eg{\isachardoublequoteclose}\isanewline
%
\isadelimtheory
\isanewline
%
\endisadelimtheory
%
\isatagtheory
\isacommand{end}\isamarkupfalse%
%
\endisatagtheory
{\isafoldtheory}%
%
\isadelimtheory
%
\endisadelimtheory
%
\end{isabellebody}%
\endinput
%:%file=Primitive_Recursion.tex%:%
%:%6=1%:%
%:%14=3%:%
%:%15=3%:%
%:%16=4%:%
%:%17=5%:%
%:%26=7%:%
%:%27=8%:%
%:%28=9%:%
%:%29=10%:%
%:%31=25%:%
%:%32=26%:%
%:%33=26%:%
%:%34=27%:%
%:%35=28%:%
%:%36=29%:%
%:%37=30%:%
%:%38=30%:%
%:%39=31%:%
%:%40=32%:%
%:%41=32%:%
%:%42=33%:%
%:%43=34%:%
%:%44=34%:%
%:%45=35%:%
%:%46=36%:%
%:%47=37%:%
%:%48=38%:%
%:%49=38%:%
%:%52=39%:%
%:%57=40%:%
